\chapter{Fluid-mediated stuff}

\section{Numerical Method}
\subsection{Governing Equations}

The fluid-structure interaction during the approach and rebound phases is governed by the coupling of the dimensionless Reynolds equation, linear elasticity, and rigid-body kinematics. 
The fluid pressure $p(\boldsymbol{x}, t)$ within the thin film is described by the transient Reynolds equation:
\begin{equation}
    \frac{\partial h}{\partial t} = \text{St} \, \nabla \cdot \left( h^3 \nabla p \right)
\end{equation}
where $\text{St}$ is the dimensionless Stokes number and $h(\boldsymbol{x}, t)$ is the dimensionless film thickness.
The gap profile is geometrically defined as the sum of the rigid-body separation $D(t)$, the static macroscopic shape of the indenter $s(\boldsymbol{x})$, and the elastic deformation $w(\boldsymbol{x}, t)$:
\begin{equation}
    h(\boldsymbol{x}, t) = D(t) + s(\boldsymbol{x}) + w(\boldsymbol{x}, t)
\end{equation}
Assuming the solid responds quasi-statically, the elastic deformation is computed using the Boussinesq integral for a linear elastic half-space:
\begin{equation}
    w(\boldsymbol{x}, t) = \frac{1}{\pi E^*} \iint_{\Omega} \frac{p(\boldsymbol{x}', t)}{|\boldsymbol{x} - \boldsymbol{x}'|} , d\boldsymbol{x}'
\end{equation}
where $E^*$ is the effective elastic modulus.
The macroscopic motion of the indenter is dictated by Newton's second law, $\dot{D}(t) = V(t)$ and $\dot{V}(t) = F(t)$, where the total hydrodynamic force $F(t)$ is obtained by integrating the pressure field over the spatial domain $\Omega$:
\begin{equation}
    F(t) = \iint_{\Omega} p(\boldsymbol{x}, t) , d\boldsymbol{x}
\end{equation}
Within the thin-film approximation of lubrication theory, the surface slope is assumed to be small ($\nabla s \ll 1$). 
Consequently, the pressure vectors are dominantly aligned with the vertical axis, and any lateral force components generating shear or torque on the indenter are considered higher-order terms and are neglected.

\subsection{Semi-Implicit Temporal Discretization}

To avoid the prohibitive computational cost of solving a fully coupled non-linear Newton-Raphson system at each time step, following Liu CITE we employ a semi-implicit backward Euler discretization. 
The rigid-body kinematics and elastic deformation are evaluated implicitly to maintain stability, while the highly non-linear fluid conductance term ($h^3$) is explicitly frozen at the known previous time step, $t^n$.
Setting a constant time step of size $\Delta t$, the time derivative of the film thickness is approximated as:
\begin{equation}
    \frac{h^{n+1} - h^n}{\Delta t} = V^{n+1} + \frac{w^{n+1} - w^n}{\Delta t}
\end{equation}
Substituting this into the Reynolds equation yields a linear system for the unknown future pressure $p^{n+1}$:
\begin{equation}
w^{n+1} - \Delta t , \text{St} , \nabla \cdot \left( (h^n)^3 \nabla p^{n+1} \right) = w^n - \Delta t V^n
\end{equation}

\subsection{Spatial Discretization}

The spatial domain is discretized using a uniform Cartesian grid with dimensions $N \times N$ and spacing $\Delta x = \Delta y$.
The Reynolds diffusion operator is discretized using a conservative Finite Volume approach.
The pressure and film thickness are evaluated at the cell centers, while the fluid fluxes are evaluated at the cell faces using staggered grid interpolation.
For a given cell $(i, j)$, the discrete divergence operator is formulated as:
\begin{equation}
    \nabla \cdot \left( h^3 \nabla p \right) \approx \frac{F_e - F_w}{\Delta x} + \frac{F_n - F_s}{\Delta y}
\end{equation}
where the face fluxes are computed utilizing a simple arithmetic average for the interface film thickness:
\begin{align}
F_e &= \left( \frac{h_{i,j} + h_{i+1,j}}{2} \right)^3 \frac{p_{i+1,j} - p_{i,j}}{\Delta x} \\
F_w &= \left( \frac{h_{i-1,j} + h_{i,j}}{2} \right)^3 \frac{p_{i,j} - p_{i-1,j}}{\Delta x} \\
F_n &= \left( \frac{h_{i,j} + h_{i,j+1}}{2} \right)^3 \frac{p_{i,j+1} - p_{i,j}}{\Delta y} \\
F_s &= \left( \frac{h_{i,j-1} + h_{i,j}}{2} \right)^3 \frac{p_{i,j} - p_{i,j-1}}{\Delta y}
\end{align}
A zero-pressure Dirichlet boundary condition, $p=0$, is enforced at the edges of the computational domain.

\subsection{Matrix-Free Linear Operator and Preconditioning}

For high-resolution 3D grids (e.g., $N=1024$), building the dense analytical Jacobian for the elastic deformation matrix would require $\mathcal{O}(N^4)$ memory (exceeding terabytes of RAM for floating point decimals), rendering standard direct solvers impossible.
Thus, we recast the semi-implicit scheme as a matrix-free linear operator $\mathcal{L}(p^{n+1}) = b$, where:
\begin{equation}
    \mathcal{L}(p) = \mathcal{K} * p - \Delta t , \text{St} \, \nabla_d \cdot \left( (h^n)^3 \nabla_d p \right)
\end{equation}
and:
\begin{equation}
    b = w^n - \Delta t \, V^{n} \, .
\end{equation}
To decouple the macroscopic kinematics from the distributed pressure field, the indenter velocity is treated explicitly during the pressure solve ($V^n$).
Once $p^{n+1}$ is computed, the hydrodynamic force is integrated, and the velocity is updated explicitly via $V^{n+1} = V^n + \Delta t \, F(p^{n+1})$.
Here, the Boussinesq convolution $\mathcal{K} * p$ is evaluated in the frequency domain using the Fast Fourier Transform (FFT). 
To prevent unphysical circular convolution artifacts inherent to the discrete Fourier transform, the pressure domain is zero-padded to size $2N \times 2N$ prior to the FFT.
This linear system is solved using the BiConjugate Gradient Stabilized method (BiCGSTAB), natively compiled via JAX for GPU execution.
BiCGSTAB is specifically chosen over GMRES due to its strict $\mathcal{O}(1)$ memory footprint, which avoids the catastrophic memory limits associated with explicitly storing large Krylov subspaces on XLA graphs.
To accelerate convergence and stabilize the Reynolds diffusion operator, we apply a diagonal Jacobi preconditioner.
The diagonal coefficient for a given cell $(i, j)$ is the sum of the local elastic self-influence and the discrete Reynolds matrix trace:
\begin{equation}
    M_{i,j} = K_{0,0} + \frac{4 \, \Delta t \, \text{St}}{\Delta x^2} (h_{i,j}^n)^3
\end{equation}
where $K_{0,0}$ is the analytical integral of the Boussinesq kernel over a single square cell $[-\Delta x/2, \Delta x/2]^2$:
\begin{equation}
    K_{0,0} = \frac{4 \Delta x}{\pi E^*} \ln(1+\sqrt{2})
\end{equation}
The linear system is preconditioned by multiplying both sides by $M^{-1}$, effectively normalizing the heterogeneous stiffness introduced by the cubic film thickness field and ensuring robust solver performance throughout the impact cycle.
Once the pressure is computed, the deformation is found by evaluating the discrete Boussinesq convolution $w^{n+1} = \mathcal{K} * p^{n+1}$ in the frequency domain, where the precomputed elastic kernel is multiplied point-wise with the transformed pressure field using the Fast Fourier Transform.